\chapter{Launching a Fund and Raising Capital}\label{ch:fm:launching-a-fund-and-raising-capital}

A new manager with a two-year track record at a Sharpe ratio of 1.5 wants a hundred million dollars. The record is barely two standard errors from
zero; the fund loses money on every dollar below its break-even size; and the seeder's term sheet asks for a fifth of the firm's revenue for as
long as the fund exists. Each of the three numbers is negotiable, and each decides whether the fund survives its first three years.

\section{The first hundred million}

\begin{definition}[Emerging manager]\label{def:fm:launching-a-fund-and-raising-capital:emerging}
An \emph{emerging manager}\index{emerging manager} is an investment manager early in its life, typically with a short track record and assets below
the size that large allocators require, often run by people who managed money at another firm before.
\end{definition}

A fund's costs arrive before its assets. People, the administrator, the auditor, lawyers, data and technology (chapter 4's service providers) cost
the chapter's manager \$2.2 million a year before a dollar is managed, plus 5 basis points of assets. Its fees are 1.5\% of assets and 20\% of gains:
on a 10\% gross year, 3.2\% of assets (\cref{lst:fm:launching-a-fund-and-raising-capital:model}, with Book 1's fee engine). The fund breaks even at
\$69.8 million of assets; with a seeder taking 20\% of the fee revenue, at \$87.6 million.

\begin{proposition}[Break-even size]\label{prop:fm:launching-a-fund-and-raising-capital:be}
With fixed costs $F$ a year, variable costs $v$ per dollar of assets, fee revenue $f(g)$ per dollar in a year of gross return $g$ and a revenue
share $s$, the fund covers its costs from assets of
\[
A^*=\frac{F}{(1-s)\,f(g)-v}.
\]
\end{proposition}

\begin{proof}
Net revenue $(1-s)f(g)A$ covers $F+vA$ exactly at $A^*$.
\end{proof}

\omcode{code/firm/fundlaunch/firm_fundlaunch.py}{43}{62}{Fees from firm.fees, the break-even size, and the t-statistic of a Sharpe ratio and the years
it needs.}\label{lst:fm:launching-a-fund-and-raising-capital:model}

\section{Seed capital and its price}

\begin{definition}[Seed investor, revenue-share agreement, capacity right]\label{def:fm:launching-a-fund-and-raising-capital:seed}
A \emph{seed investor}\index{seed investor} provides a new manager's first large allocation, usually locked up for several years, in exchange for
economics beyond the fund's returns. A \emph{revenue-share agreement}\index{revenue-share agreement} gives the seeder a share of the manager's gross
fee revenue, sometimes for the life of the business and sometimes until a buy-back. A \emph{capacity right}\index{capacity right} entitles an
investor to add to its allocation later, even when the fund is closed to others.
\end{definition}

\begin{definition}[Founders' share class]\label{def:fm:launching-a-fund-and-raising-capital:founders}
A \emph{founders' share class}\index{founders' share class} offers the earliest investors lower fees, for example a lower management fee or
performance fee, for as long as they stay invested, often closed once the fund reaches a stated size.
\end{definition}

The seeder's share is expensive: it raises the break-even by \$17.8 million and takes a fifth of every future fee. It buys time. The chapter's manager
either launches with \$50 million from a seeder that takes 20\% of revenue, or with \$15 million from friends and family and no seeder; everything
else is the same (\cref{tab:fm:launching-a-fund-and-raising-capital:cases}).

\section{What a track record proves}

A Sharpe ratio estimated from $T$ years of returns has a standard error of about $1/\sqrt T$ (Book 4), so its $t$-statistic is $S\sqrt T$. A record
at a Sharpe ratio of 1.5 over two years has $t=2.12$; a Sharpe ratio of 1 needs four years to reach $t=2$, and 0.5 needs sixteen. Investors who
wait for statistical evidence therefore wait years, and the fund must survive them. Book 4's deflated Sharpe ratio sharpens the point: a manager who
tried several strategies before launching one has a weaker record than its $t$-statistic says.

\begin{definition}[Track record portability]\label{def:fm:launching-a-fund-and-raising-capital:port}
\emph{Track record portability}\index{track record portability} is the extent to which a manager may present performance achieved at a previous firm
as its own: in the United States it is allowed in advertisements only under the conditions of the SEC's marketing rule
(\cref{dat:fm:launching-a-fund-and-raising-capital:marketing}).
\end{definition}

\begin{dated}{2026-09}{Advertising performance in the United States}\label{dat:fm:launching-a-fund-and-raising-capital:marketing}
The SEC's marketing rule (17 CFR 275.206(4)-1, adopted in 2021) forbids an adviser to advertise gross performance without net performance of at
least equal prominence; requires one-, five- and ten-year periods for performance other than private funds'; and allows predecessor performance only
if the people primarily responsible for it manage accounts at the advertising adviser, the accounts are sufficiently similar, all substantially
similar accounts are included unless excluding them would not materially raise the results, and the advertisement discloses that the performance
was achieved at another entity. This box summarises the rule's text and is not legal advice.
\end{dated}

Aggarwal and Jorion's study of emerging hedge funds and managers (2010) is the reference on whether new funds perform differently from established
ones; for the chapter the point is narrower: whatever a new fund's true skill, its investors can see little of it for years.

\section{Tutorial: the first hundred million}\label{tut:fm:launching-a-fund-and-raising-capital:tut}

\begin{tutorial}
\textbf{Goal.} Simulate a fund's first five years with and without a seed deal: its assets, the probability of reaching \$100 million, and the value
of the deal to the manager and the seeder. \textbf{End state:} the fan chart (\cref{fig:fm:launching-a-fund-and-raising-capital:fan}) and the
probability curves (\cref{fig:fm:launching-a-fund-and-raising-capital:reach}).
\begin{tutsteps}
\item \textbf{The strategy.} A true Sharpe ratio of 1 at 10\% volatility; monthly returns drawn independently (illustrative).
\item \textbf{The investors.} Each month assets grow by an inflow of up to 4\% of assets, scaled by a logistic function of the track record's
$t$-statistic (half confidence at $t=2$, none in the first six months), and shrink by 10\% in months when the fund is more than 10\% below its peak.
\item \textbf{The economics.} \texttt{firm.fundlaunch.simulate} gives 2\,000 paths of assets and net fees; \texttt{manager\_value} discounts fees
less costs at 8\%; \texttt{seed\_value} values the seeder's revenue share.
\item \textbf{The comparison.} \texttt{fm\_launch.summary()} for the seeded and the unseeded launch.
\end{tutsteps}
\end{tutorial}

\begin{table}[ht]
\centering\small
\begin{tabular}{lrr}
\toprule
 & seeded, 20\% share & unseeded, \$15 million \\
\midrule
break-even assets (\$ million) & 87.6 & 69.8 \\
probability of \$100 million within 3 years & 43.2\% & 0.25\% \\
median month of reaching it (paths that do, in 5 years) & 32 & 51 \\
assets at 3 years, 10th / 50th / 90th percentile & 30.7 / 88.9 / 186.0 & 9.2 / 26.7 / 55.8 \\
below break-even after 5 years & 32.3\% & 64.5\% \\
manager's value, 5 years (\$ million) & 7.52 & $-2.85$ \\
seeder's revenue share, 5 years (\$ million) & 4.21 & -- \\
\bottomrule
\end{tabular}
\caption{A fund's first five years with and without a seeder (2\,000 simulated paths each; values discounted at 8\%). Data:
\texttt{fm\_launch.summary}.}\label{tab:fm:launching-a-fund-and-raising-capital:cases}
\end{table}

\begin{omfigure}
\begin{tikzpicture}
\begin{axis}[width=10.5cm,height=5.4cm,xlabel={\small month},ylabel={\small assets (\$ million)},xmin=1,xmax=60,ymin=0,ymax=600,
  tick label style={font=\footnotesize},grid=major,grid style={gray!20},table/col sep=comma,
  legend cell align=left,legend style={font=\scriptsize,at={(0.5,-0.3)},anchor=north,/tikz/every even column/.append style={column sep=8pt}},legend columns=3]
\addplot[name path=p90,draw=none,forget plot] table[x=month,y=p90]{figdata/desk/25-launching-a-fund-and-raising-capital/fan.csv};
\addplot[name path=p10,draw=none,forget plot] table[x=month,y=p10]{figdata/desk/25-launching-a-fund-and-raising-capital/fan.csv};
\addplot[omDef!15] fill between[of=p90 and p10];
\addplot[name path=p75,draw=none,forget plot] table[x=month,y=p75]{figdata/desk/25-launching-a-fund-and-raising-capital/fan.csv};
\addplot[name path=p25,draw=none,forget plot] table[x=month,y=p25]{figdata/desk/25-launching-a-fund-and-raising-capital/fan.csv};
\addplot[omDef!35] fill between[of=p75 and p25];
\addplot[omDef,thick] table[x=month,y=p50]{figdata/desk/25-launching-a-fund-and-raising-capital/fan.csv};
\draw[omProp,dashed,thick] (axis cs:1,87.6) -- (axis cs:60,87.6);
\legend{10th--90th,25th--75th,median}
\end{axis}
\end{tikzpicture}

\omcaption{Assets of the seeded fund over its first five years: percentiles of 2\,000 simulated paths, and the break-even size with the seeder's
revenue share (dashed, \$87.6 million). A third of the paths are still below it after five years. Data: \texttt{fm\_launch.fan}.}\label{fig:fm:launching-a-fund-and-raising-capital:fan}
\end{omfigure}

The seed changes everything in the model. With \$50 million the fund has a 43\% chance of reaching \$100 million within three years, typically in
month 32; with \$15 million the chance is a quarter of one per cent, since the inflows that a track record attracts are proportional to the assets
already there, and a small fund grows slowly even when its record is good. Over five years the seeded manager's business is worth \$7.52 million
after the seeder's \$4.21 million; the unseeded manager's loses \$2.85 million. The seed is worth \$10.37 million to the manager, and the revenue share
costs it \$4.21 million of that: a price worth paying in the model, and a large one
(\cref{fig:fm:launching-a-fund-and-raising-capital:values}).

\begin{omfigure}
\begin{tikzpicture}
\begin{axis}[width=10.5cm,height=5cm,xlabel={\small months since launch},ylabel={\small probability of \$100 million (\%)},xmin=1,xmax=60,ymin=0,ymax=80,
  tick label style={font=\footnotesize},grid=major,grid style={gray!20},table/col sep=comma,
  legend cell align=left,legend style={font=\scriptsize,at={(0.5,-0.3)},anchor=north,/tikz/every even column/.append style={column sep=8pt}},legend columns=2]
\addplot[omDef,thick] table[x=month,y=seeded]{figdata/desk/25-launching-a-fund-and-raising-capital/reach.csv};
\addplot[omProp,thick,dashed] table[x=month,y=unseeded]{figdata/desk/25-launching-a-fund-and-raising-capital/reach.csv};
\draw[gray,dotted] (axis cs:36,0) -- (axis cs:36,80);
\legend{{seeded, \$50 million},{unseeded, \$15 million}}
\end{axis}
\end{tikzpicture}

\omcaption{The probability that the fund has reached \$100 million by each month, with and without the seed. Data: \texttt{fm\_launch.reach\_curve}.}\label{fig:fm:launching-a-fund-and-raising-capital:reach}
\end{omfigure}

\begin{omfigure}
\begin{tikzpicture}
\begin{axis}[width=9cm,height=5cm,ybar,bar width=16pt,ylabel={\small present value, 5 years (\$ million)},ymin=-4,ymax=10,xtick={0,1},
  xticklabels={{seeded, 20\% share},{unseeded, \$15 million}},xmin=-0.6,xmax=1.6,tick label style={font=\footnotesize},grid=major,
  grid style={gray!20},table/col sep=comma,nodes near coords,nodes near coords style={font=\scriptsize},
  legend cell align=left,legend style={font=\scriptsize,at={(0.5,-0.25)},anchor=north,/tikz/every even column/.append style={column sep=8pt}},legend columns=2]
\addplot[fill=omDef!55,draw=omDef] table[x=pos,y=manager]{figdata/desk/25-launching-a-fund-and-raising-capital/values.csv};
\addplot[fill=omProp!55,draw=omProp] table[x=pos,y=seeder]{figdata/desk/25-launching-a-fund-and-raising-capital/values.csv};
\draw[black,thin] (axis cs:-0.6,0) -- (axis cs:1.6,0);
\legend{manager,seeder's revenue share}
\end{axis}
\end{tikzpicture}

\omcaption{The seed deal's value over five years to the manager and the seeder, against launching alone with \$15 million. Data:
\texttt{fm\_launch.summary}.}\label{fig:fm:launching-a-fund-and-raising-capital:values}
\end{omfigure}

\section{Due diligence}

\begin{definition}[Operational due diligence, due diligence questionnaire]\label{def:fm:launching-a-fund-and-raising-capital:odd}
\emph{Operational due diligence}\index{operational due diligence} is an investor's review of a manager's non-investment functions: governance,
valuation, trade processing, cash controls, service providers, compliance, technology and business continuity. A \emph{due diligence
questionnaire}\index{due diligence questionnaire} is the standard document of questions a manager answers for it, usually in an industry
association's format, with the supporting policies attached.
\end{definition}

\omterm{def:fm:launching-a-fund-and-raising-capital:odd}{Operational due diligence} fails funds that pass the investment review: a single person able to move cash, marks set by the trader (chapter 15), an
administrator that only records what the manager tells it. The chapters of this book are the answers a new manager needs: the risk function
(chapter 12), operations (chapter 13), compliance (chapter 16), the legal documents (chapter 18), crisis plans (chapter 27).

\begin{method}[Raising the first hundred million]\label{met:fm:launching-a-fund-and-raising-capital:raise}
\begin{enumerate}
\item Compute the break-even size and the months of costs the firm can fund before reaching it, on pessimistic inflows.
\item State what the track record proves, with its $t$-statistic and the number of strategies tried, and what it cannot.
\item Negotiate the seed on the revenue share's size, duration and buy-back, the lock-up, \omterm{def:fm:launching-a-fund-and-raising-capital:seed}{capacity rights} and information rights, valued on the
firm's own model of its growth.
\item Prepare \omterm{def:fm:launching-a-fund-and-raising-capital:odd}{operational due diligence} before the first meeting: policies, service providers, controls, a completed questionnaire.
\item Offer a founders' class to early investors, closed at a stated size.
\end{enumerate}
\end{method}

\section{Raising: investors, consultants and the pitch}

An allocator that will not own more than a small share of any fund cannot invest much in a small one, and one that waits for a statistically
convincing record waits years (the previous section); the consultants many allocators use screen the same way. The early money therefore comes
from the manager's network, seeders, family offices, funds of \omterm{def:fm:launching-a-fund-and-raising-capital:emerging}{emerging managers} and platforms (chapter 3), each with its own price. The pitch is the track record, the strategy's capacity and the operational
answers; the negotiation is over fees, liquidity and the \omterm{def:fm:the-asset-manager-and-the-fund:side}{side letters} (chapter 4) that the largest early investors ask for, subject to the
\omterm{def:fm:the-asset-manager-and-the-fund:side}{most-favoured-nation clauses} they will later trigger.

\section{Build: the launch model}\label{bld:fm:launching-a-fund-and-raising-capital:fundlaunch}

\begin{build}
\bfield{Purpose} A fund launch's economics: costs, fees net of the revenue share, break-even size, inflows driven by the track record, the
probability of reaching a size, and the seed deal's value to both sides.
\bfield{Interface} \texttt{firm.fundlaunch}: \texttt{Launch}, \texttt{annual\_fees} (on \texttt{firm.fees}), \texttt{breakeven\_aum}, \texttt{tstat},
\texttt{years\_for\_t}, \texttt{simulate}, \texttt{prob\_reach}, \texttt{first\_month}, \texttt{manager\_value}, \texttt{seed\_value}.
\bfield{Rules} Fees are computed by Book 1's engine; the revenue share is taken from gross fees; no inflows before six months of record; assets
never fall below zero.
\bfield{Acceptance tests} \texttt{code/firm/fundlaunch/tests/}: fees and break-even on hand numbers; path shapes; the helpers on small arrays.
\bfield{Stretch} A buy-back of the revenue share; founders' class fees; the seeder's own return on its capital; serial correlation of returns.
\end{build}

\begin{omsources}
\item 17 CFR 275.206(4)-1 (the SEC marketing rule).
\item V.~Aggarwal and P.~Jorion, ``The performance of emerging hedge funds and managers'', \emph{Journal of Financial Economics}, 2010.
\end{omsources}

\section{Exercises}

\begin{exercise}[$\star$]\label{exo:fm:launching-a-fund-and-raising-capital:1}
Compute the fee revenue per dollar in a 10\% gross year at 1.5 and 20, and the break-even size with and without a 20\% revenue share.
\end{exercise}

\begin{exercise}[$\star$]\label{exo:fm:launching-a-fund-and-raising-capital:2}
How many years does a Sharpe ratio of 1 need to be two standard errors from zero? A Sharpe ratio of 0.5?
\end{exercise}

\begin{exercise}[$\star$]\label{exo:fm:launching-a-fund-and-raising-capital:3}
Which conditions must a manager meet to advertise performance achieved at its previous firm in the United States?
\end{exercise}

\begin{exercise}[$\star\star$]\label{exo:fm:launching-a-fund-and-raising-capital:4}
Why does the unseeded fund almost never reach \$100 million within three years, although its strategy is the same?
\end{exercise}

\begin{exercise}[$\star\star$]\label{exo:fm:launching-a-fund-and-raising-capital:5}
What is the seed worth to the manager over five years, and what does the revenue share cost it?
\end{exercise}

\begin{exercise}[$\star\star$]\label{exo:fm:launching-a-fund-and-raising-capital:6}
Why do \omterm{def:fm:launching-a-fund-and-raising-capital:odd}{operational due diligence} failures stop funds that pass the investment review? Give three examples.
\end{exercise}

\begin{exercise}[$\star\star\star$]\label{exo:fm:launching-a-fund-and-raising-capital:7}
\emph{Coding.} Rerun the seeded launch with a revenue share of 10\%. What happens to the break-even size and the manager's value?
\end{exercise}

\begin{exercise}[$\star\star\star$]\label{exo:fm:launching-a-fund-and-raising-capital:8}
\emph{Find the flaw.} ``Our two-year Sharpe ratio of 1.5 proves our skill; investors will come.''
\end{exercise}

\section{Problem: The First Hundred Million}

\begin{problem}[{Weekend problem --- the first hundred million}]\label{pb:fm:launching-a-fund-and-raising-capital:1}
A manager leaving a large firm must decide whether to accept a seeder's offer or launch alone.

\medskip\noindent\textbf{Part I --- The economics.}
\begin{enumerate}
\item Define an \omterm{def:fm:launching-a-fund-and-raising-capital:emerging}{emerging manager}.
\item State and prove \cref{prop:fm:launching-a-fund-and-raising-capital:be}.
\item Give the break-even size with and without the revenue share.
\item Define a \omterm{def:fm:launching-a-fund-and-raising-capital:seed}{seed investor}, a \omterm{def:fm:launching-a-fund-and-raising-capital:seed}{revenue-share agreement}, a \omterm{def:fm:launching-a-fund-and-raising-capital:seed}{capacity right} and a \omterm{def:fm:launching-a-fund-and-raising-capital:founders}{founders' share class}.
\end{enumerate}

\medskip\noindent\textbf{Part II --- The record.}
\begin{enumerate}[resume]
\item What is the standard error of a Sharpe ratio, and the $t$-statistic of 1.5 over two years?
\item How long does a Sharpe ratio of 1 take to reach $t=2$?
\item Define \omterm{def:fm:launching-a-fund-and-raising-capital:port}{track record portability} and state the US rule.
\item How does the number of strategies tried change what the record proves?
\end{enumerate}

\medskip\noindent\textbf{Part III --- The launch.}
\begin{enumerate}[resume]
\item Describe the inflow model.
\item Give the probability of reaching \$100 million within three years, seeded and unseeded.
\item Give the assets at three years and the share below break-even after five.
\item Give the manager's and the seeder's values.
\item What does the model leave out?
\end{enumerate}

\medskip\noindent\textbf{Part IV --- The decision.}
\begin{enumerate}[resume]
\item Define \omterm{def:fm:launching-a-fund-and-raising-capital:odd}{operational due diligence} and a \omterm{def:fm:launching-a-fund-and-raising-capital:odd}{due diligence questionnaire}.
\item What would you negotiate in the seed deal?
\item Would you accept the seeder's offer?
\item Who else could provide early capital, and at what price?
\item What would you prepare before the first investor meeting?
\item State the \emph{named result}: the break-even AUM and the months to a hundred million dollars under a 20 per cent revenue share, and the
track-record length at which a Sharpe ratio of 1 is two standard errors from zero.
\item In two sentences, write the launch plan.
\end{enumerate}
\end{problem}

\section{Interview questions}

\begin{interviewq}[$\star$ \iqroles{researcher}]\label{iq:fm:launching-a-fund-and-raising-capital:1}
What is the standard error of an annual Sharpe ratio estimated from five years of data?
\end{interviewq}

\begin{interviewq}[$\star$ \iqroles{trader}]\label{iq:fm:launching-a-fund-and-raising-capital:2}
Why do new funds offer founders' share classes?
\end{interviewq}

\begin{interviewq}[$\star\star$ \iqroles{researcher}]\label{iq:fm:launching-a-fund-and-raising-capital:3}
A fund's costs are \$2 million a year and it earns 2\% of assets in fees. What is its break-even size, and how does a 25\% revenue share change it?
\end{interviewq}

\begin{interviewq}[$\star\star$ \iqroles{risk}]\label{iq:fm:launching-a-fund-and-raising-capital:4}
What would you look for in \omterm{def:fm:launching-a-fund-and-raising-capital:odd}{operational due diligence} of a new manager?
\end{interviewq}

\begin{interviewq}[$\star\star$ \iqroles{trader, researcher}]\label{iq:fm:launching-a-fund-and-raising-capital:5}
Can you show investors the track record you built at your previous firm?
\end{interviewq}

\begin{interviewq}[$\star\star\star$ \iqroles{researcher}]\label{iq:fm:launching-a-fund-and-raising-capital:6}
How would you value a seeder's revenue share, and what terms would you trade against it?
\end{interviewq}
